\chapter{2012年江苏卷}

\section{填空题}

\begin{problem}
已知集合 \(A = \{1, 2, 4\}\)，\(B = \{2, 4, 6\}\)，则 \(A \cup B =\fillinblank{}\).
\end{problem}

\begin{answer}
\(\{1,2,4,6\}\)
\end{answer}

\begin{solution}
并集包含属于集合 \(A\) 或集合 \(B\) 的所有元素，因此把两个集合中的元素合并并去掉重复元素，得到
\[
A\cup B=\{1,2,4,6\}
\]
所以填 \(\{1,2,4,6\}\).
\end{solution}

\begin{problem}
某学校高一、高二、高三年级的学生人数之比是 \(3:3:4\). 现用分层抽样的方法从该校高中三个年级的学生中抽取容量为 \(50\) 的样本. 则应从高二年级抽取\fillinblank{} 名学生.
\end{problem}

\begin{answer}
15
\end{answer}

\begin{solution}
三个年级的学生人数之比为 \(3:3:4\)，所以高二年级人数占三个年级总人数的
\[
\frac{3}{3+3+4}=\frac{3}{10}
\]
分层抽样保持各层的比例不变，故应抽取的高二年级学生数为
\[
50\times\frac{3}{10}=15
\]
所以填 \(15\).
\end{solution}

\begin{problem}
设 \(a\)，\(b \in \R\)，\(a + b\i = \frac{11 - 7\i}{1 - 2\i}\)（\(\i\) 为虚数单位），则 \(a + b\) 的值为\fillinblank{}.
\end{problem}

\begin{answer}
8
\end{answer}

\begin{solution}
将分式的分子、分母同乘以分母的共轭复数 \(1+2\i\)，得
\[
\frac{11-7\i}{1-2\i}=\frac{(11-7\i)(1+2\i)}{(1-2\i)(1+2\i)}=\frac{25+15\i}{5}=5+3\i
\]
比较实部和虚部，得到 \(a=5\)，\(b=3\)，所以
\[
a+b=8
\]
\end{solution}

\begin{problem}
如图是一个算法流程图，则输出的 \(k\) 的值是\fillinblank{}.

\bitmapfigure[max width=0.46\linewidth]{img/2012/jiangsu/q4_fig1.png}
\end{problem}

\begin{answer}
5
\end{answer}

\begin{solution}
流程开始时 \(k=1\). 当 \(k=1\) 时，\(k^2-5k+4=0\)，不满足大于 \(0\)，所以按流程令 \(k=2\). 当 \(k=2\) 时，\(k^2-5k+4=-2\)，仍令 \(k=3\). 当 \(k=3\) 时，该式为 \(-2\)，令 \(k=4\). 当 \(k=4\) 时，该式为 \(0\)，令 \(k=5\). 当 \(k=5\) 时，该式为 \(4>0\)，流程输出 \(k\). 因此输出值为 \(5\).
\end{solution}

\begin{problem}
函数 \(f(x)=\sqrt{1-2\log_{6}x}\) 的定义域为\fillinblank{}.
\end{problem}

\begin{answer}
\(\left(0,\sqrt{6}\right]\)
\end{answer}

\begin{solution}
要使 \(f(x)\) 有意义，首先必须有 \(x>0\)，并且根式内的数非负，即
\[
1-2\log_{6}x\geqslant0
\]
由于底数 \(6>1\)，上式等价于
\[
\log_{6}x\leqslant\frac{1}{2}
\]
从而 \(x\leqslant6^{1/2}=\sqrt{6}\). 合并两个条件，函数的定义域为
\[
\left(0,\sqrt{6}\right]
\]
\end{solution}

\begin{problem}
现有 \(10\) 个数，它们能构成一个以 \(1\) 为首项，\(-3\) 为公比的等比数列，若从这 \(10\) 个数中随机抽取一个数，则它小于 \(8\) 的概率是\fillinblank{}.
\end{problem}

\begin{answer}
\(\frac{3}{5}\)
\end{answer}

\begin{solution}
这 \(10\) 个数依次为
\(1\)，\(-3\)，\(9\)，\(-27\)，\(81\)，\(-243\)，\(729\)，\(-2187\)，\(6561\)，\(-19683\).
其中第一个数 \(1<8\)，所有偶数项均为负数，也都小于 \(8\)，而从第三项开始的奇数项分别为 \(9\)，\(81\)，\(729\)，\(6561\)，均不小于 \(8\). 因此小于 \(8\) 的数共有 \(1+5=6\) 个，所求概率为
\[
\frac{6}{10}=\frac{3}{5}
\]
\end{solution}

\begin{problem}
如图，在长方体 \(ABCD-A_{1}B_{1}C_{1}D_{1}\) 中，\(AB = AD = 3 \mathrm{~cm}\)，\(AA_{1} = 2 \mathrm{~cm}\)，则四棱锥 \(A-BB_{1}D_{1}D\) 的体积为\fillinblank{} \(\mathrm{cm}^{3}\).

\bitmapfigure[max width=0.42\linewidth]{img/2012/jiangsu/q7_fig1.png}
\end{problem}

\begin{answer}
6
\end{answer}

\begin{solution}
四棱锥 \(A-BB_{1}D_{1}D\) 的底面 \(BB_{1}D_{1}D\) 是矩形，且
\[
BD=\sqrt{AB^{2}+AD^{2}}=3\sqrt{2}
\]
且 \(BB_{1}=AA_{1}=2\)，
所以底面积为
\[
S_{BB_{1}D_{1}D}=BD\cdot BB_{1}=6\sqrt{2}
\]
在正方形 \(ABCD\) 中，三角形 \(ABD\) 的面积为 \(9/2\)，故点 \(A\) 到直线 \(BD\) 的距离为
\[
\frac{2\times(9/2)}{3\sqrt{2}}=\frac{3}{\sqrt{2}}
\]
底面所在平面包含直线 \(BD\) 并垂直于底面 \(ABCD\)，所以点 \(A\) 到平面 \(BB_{1}D_{1}D\) 的距离也是 \(3/\sqrt{2}\). 因而四棱锥的体积为
\[
V=\frac{1}{3}\times6\sqrt{2}\times\frac{3}{\sqrt{2}}=6
\]
所以填 \(6\).
\end{solution}

\begin{problem}
在平面直角坐标系 \(xOy\) 中，若双曲线 \(\frac{x^2}{m} - \frac{y^2}{m^2 + 4} = 1\) 的离心率为 \(\sqrt{5}\)，则 \(m\) 的值为\fillinblank{}.
\end{problem}

\begin{answer}
2
\end{answer}

\begin{solution}
双曲线的标准方程中 \(a^{2}=m\)，因此 \(m>0\)，且
\[
c^{2}=a^{2}+b^{2}=m+m^{2}+4
\]
由离心率 \(e=\sqrt{5}\)，有
\[
e^{2}=\frac{c^{2}}{a^{2}}=\frac{m^{2}+m+4}{m}=5
\]
于是
\[
m+1+\frac{4}{m}=5
\]
即
\[
m^{2}-4m+4=0
\]
所以 \((m-2)^{2}=0\)，得到 \(m=2\).
\end{solution}

\begin{problem}
如图，在矩形 \(ABCD\) 中，\(AB = \sqrt{2}\)，\(BC = 2\)，点 \(E\) 为 \(BC\) 的中点，点 \(F\) 在边 \(CD\) 上，若 \(\overrightarrow{AB} \cdot \overrightarrow{AF} = \sqrt{2}\)，则 \(\overrightarrow{AE} \cdot \overrightarrow{BF}\) 的值是\fillinblank{}.

\bitmapfigure[max width=0.34\linewidth]{img/2012/jiangsu/q9_fig1.png}
\end{problem}

\begin{answer}
\(\sqrt{2}\)
\end{answer}

\begin{solution}
建立直角坐标系，取 \(A=(0,0)\)，\(B=(\sqrt{2},0)\)，\(C=(\sqrt{2},2)\)，\(D=(0,2)\). 因为 \(E\) 是 \(BC\) 的中点，所以 \(E=(\sqrt{2},1)\). 设 \(F=(t,2)\)，其中 \(0\leqslant t\leqslant\sqrt{2}\). 则
\[
\overrightarrow{AB}=(\sqrt{2},0)
\]
且 \(\overrightarrow{AF}=(t,2)\)，
由题设
\[
\overrightarrow{AB}\cdot\overrightarrow{AF}=\sqrt{2}t=\sqrt{2}
\]
得 \(t=1\)，所以 \(F=(1,2)\). 此时
\[
\overrightarrow{AE}=(\sqrt{2},1)
\]
且 \(\overrightarrow{BF}=(1-\sqrt{2},2)\)，
从而
\[
\overrightarrow{AE}\cdot\overrightarrow{BF}=\sqrt{2}(1-\sqrt{2})+2=\sqrt{2}
\]
\end{solution}

\begin{problem}
设 \(f(x)\) 是定义在 \(\R\) 上且周期为 \(2\) 的函数，且在区间 \([-1,1]\) 上， \(f(x)=\begin{cases}
ax+1, & -1 \leqslant x < 0,\\
\frac{bx+2}{x+1}, & 0 \leqslant x \leqslant 1,
\end{cases}\)
其中 \(a\)，\(b \in \R\). 若 \(f\left(\frac{1}{2}\right)=f\left(\frac{3}{2}\right)\)，则 \(a + 3b\) 的值为\fillinblank{}.
\end{problem}

\begin{answer}
\(-10\)
\end{answer}

\begin{solution}
由周期为 \(2\)，且 \(\frac{3}{2}-\left(-\frac{1}{2}\right)=2\)，有
\[
f\left(\frac{3}{2}\right)=f\left(-\frac{1}{2}\right)
\]
根据分段表达式，条件 \(f\left(\frac{1}{2}\right)=f\left(\frac{3}{2}\right)\) 化为
\[
\frac{b+4}{3}=1-\frac{a}{2}
\]
即
\[
3a+2b=-2
\]
又因为函数周期为 \(2\)，所以 \(f(-1)=f(1)\). 代入分段表达式得
\[
1-a=\frac{b+2}{2}
\]
即 \(b=-2a\). 联立两个方程，得 \(a=2\)，\(b=-4\). 因此
\[
a+3b=2+3(-4)=-10
\]
\end{solution}

\begin{problem}
设 \(\alpha\) 为锐角，若 \(\cos\left(\alpha + \frac{\pi}{6}\right)=\frac{4}{5}\)，则 \(\sin\left(2\alpha + \frac{\pi}{12}\right)\) 的值为\fillinblank{}.
\end{problem}

\begin{answer}
\(\frac{17\sqrt{2}}{50}\)
\end{answer}

\begin{solution}
令 \(\beta=\alpha+\frac{\pi}{6}\). 由于 \(\alpha\) 为锐角，且
\(\frac{\pi}{6}<\beta<\frac{2\pi}{3}\)，\(\cos\beta=\frac45>0\)
所以在上述范围内 \(\beta\) 必为锐角，从而
\[
\sin\beta=\frac35
\]
由 \(\alpha=\beta-\frac{\pi}{6}\)，有
\[
2\alpha+\frac{\pi}{12}=2\beta-\frac{\pi}{4}
\]
又
\[
\sin2\beta=2\sin\beta\cos\beta=\frac{24}{25}
\]
又 \(\cos2\beta=\cos^{2}\beta-\sin^{2}\beta=\frac{7}{25}\)，
因此
\[
\sin\left(2\beta-\frac{\pi}{4}\right)=\frac{\sin2\beta-\cos2\beta}{\sqrt{2}}=\frac{17}{25\sqrt{2}}=\frac{17\sqrt{2}}{50}
\]
\end{solution}

\begin{problem}
在平面直角坐标系 \(xOy\) 中，圆 \(C\) 的方程为 \(x^2 + y^2 - 8x + 15 = 0\)，若直线 \(y = kx - 2\) 上至少存在一点，使得以该点为圆心，\(1\) 为半径的圆与圆 \(C\) 有公共点，则 \(k\) 的最大值是\fillinblank{}.
\end{problem}

\begin{answer}
\(\frac{4}{3}\)
\end{answer}

\begin{solution}
将圆 \(C\) 的方程配方，得
\[
(x-4)^{2}+y^{2}=1
\]
所以圆 \(C\) 的圆心为 \(O_{C}=(4,0)\)，半径为 \(1\). 设直线 \(y=kx-2\) 上的点为 \(Q\). 以 \(Q\) 为圆心、半径为 \(1\) 的圆与 \(C\) 有公共点，当且仅当两圆心的距离不超过两个半径之和，即
\[
QO_{C}\leqslant2
\]
因此，直线 \(y=kx-2\) 到点 \(O_{C}\) 的距离不超过 \(2\). 由点到直线的距离公式，得到
\[
\frac{|4k-2|}{\sqrt{k^{2}+1}}\leqslant2
\]
两边平方并化简，得
\[
(4k-2)^{2}\leqslant4(k^{2}+1)
\]
即
\[
4k(3k-4)\leqslant0
\]
所以 \(0\leqslant k\leqslant\frac43\)，故 \(k\) 的最大值为 \(\frac43\).
\end{solution}

\begin{problem}
已知函数 \(f(x)=x^2+ax+b\)（\(a\)，\(b\in\R\)）的值域为 \([0,+\infty)\)，若关于 \(x\) 的不等式 \(f(x)<c\) 的解集为 \((m,m+6)\)，则实数 \(c\) 的值为\fillinblank{}.
\end{problem}

\begin{answer}
9
\end{answer}

\begin{solution}
函数 \(f(x)=x^{2}+ax+b\) 的图象开口向上，且值域为 \([0,+\infty)\)，因此其最小值为 \(0\). 配方得
\[
f(x)=\left(x+\frac a2\right)^{2}+b-\frac{a^{2}}{4}
\]
所以顶点的纵坐标为 \(b-a^{2}/4=0\). 题设给出的解集非空，而 \(f(x)\) 的最小值为 \(0\)，故必须有 \(c>0\). 此时不等式 \(f(x)<c\) 的解集为
\[
\left(-\frac a2-\sqrt{c},-\frac a2+\sqrt{c}\right)
\]
其长度为 \(2\sqrt{c}\). 题设解集 \((m,m+6)\) 的长度为 \(6\)，故
\[
2\sqrt{c}=6
\]
从而 \(c=9\).
\end{solution}

\begin{problem}
已知正数 \(a\)，\(b\)，\(c\) 满足 \(5c-3a \leqslant b \leqslant 4c-a\)，\(c \ln b \geqslant a + c \ln c\)，则 \(\frac{b}{a}\) 的取值范围是\fillinblank{}.
\end{problem}

\begin{answer}
\(\left[\e,7\right]\)
\end{answer}

\begin{solution}
令 \(r=\frac ac>0\)，\(t=\frac bc>0\).
题设的两个线性不等式化为
\[
5-3r\leqslant t\leqslant4-r
\]
而对数不等式化为
\[
\ln t\geqslant r
\]

由 \(5-3r\leqslant4-r\)，得 \(r\geqslant1/2\). 又由 \(t\leqslant4-r\)，有
\[
\frac ba=\frac tr\leqslant\frac{4-r}{r}=\frac4r-1\leqslant7
\]

另一方面，\(\ln t\geqslant r>0\)，所以 \(t>1\)，且 \(r\leqslant\ln t\). 令 \(g(t)=\frac{t}{\ln t}\)，其中 \(t>1\).
则
\[
g'(t)=\frac{\ln t-1}{(\ln t)^2}
\]
当 \(1<t<\e\) 时，\(g'(t)<0\)；当 \(t>\e\) 时，\(g'(t)>0\)，所以 \(g(t)\) 在 \(t=\e\) 处取得最小值 \(g(\e)=\e\). 因此
\[
\frac ba=\frac tr\geqslant\frac{t}{\ln t}\geqslant\e
\]
于是
\[
\e\leqslant\frac ba\leqslant7
\]

下界可以取得：取 \(r=1\)，\(t=\e\)，则 \(5-3r=2\leqslant\e\leqslant3=4-r\)，且 \(\ln t=r\)，所以 \(b/a=\e\). 上界也可以取得：取 \(r=1/2\)，\(t=7/2\)，则两个线性不等式均取等号，且 \(\ln(7/2)>1/2\)，所以 \(b/a=7\).

再说明中间值均可取得. 令 \(s\in[0,1]\)，\(r_s=1-\frac{s}{2}\)，\(t_s=\e+s\left(\frac{7}{2}-\e\right)\). 两个线性不等式对 \((r_s,t_s)\) 成立，因为它们对端点 \((1,\e)\) 和 \((1/2,7/2)\) 成立. 又由 \(\ln t\) 的凹性，
\[
\ln t_s\geqslant(1-s)\ln\e+s\ln\frac{7}{2}=1+s\left(\ln\frac{7}{2}-1\right)\geqslant1-\frac{s}{2}=r_s
\]
所以 \((r_s,t_s)\) 始终满足全部条件. 比值 \(t_s/r_s\) 关于 \(s\) 连续，且在 \(s=0\) 和 \(s=1\) 时分别为 \(\e\) 和 \(7\)，由介值定理可取得两者之间的每一个值. 因而所求取值范围为
\[
\left[\e,7\right]
\]
\end{solution}

\section{解答题}

\begin{problem}
在 \(\triangle ABC\) 中，已知 \(\overrightarrow{AB} \cdot \overrightarrow{AC} = 3\overrightarrow{BA} \cdot \overrightarrow{BC}\).

\begin{enumerate}
\item 求证：\(\tan B = 3\tan A\).
\item 若 \(\cos C = \frac{\sqrt{5}}{5}\)，求 \(A\) 的值.
\end{enumerate}
\end{problem}

\begin{solution}
\begin{enumerate}
\item 记三角形 \(ABC\) 的边长分别为 \(a=BC\)，\(b=CA\)，\(c=AB\). 由向量数量积的定义，题设等式化为

\[
bc\cos A=3ac\cos B
\]

若 \(\cos A=0\)，则上式左边为 \(0\)，从而 \(\cos B=0\)，这与三角形内角和为 \(\pi\) 矛盾；若 \(\cos B=0\)，则上式右边为 \(0\)，从而 \(\cos A=0\)，同样与三角形内角和为 \(\pi\) 矛盾. 所以 \(\cos A\neq0\)，\(\cos B\neq0\). 将上式除以 \(ac\)，并由正弦定理 \(b/a=\sin B/\sin A\)，得

\[
\frac{b}{a}\cos A=3\cos B
\]

于是 \(\frac{\sin B}{\sin A}\cos A=3\cos B\). 两边同除以 \(\cos B\)，再乘以 \(\sin A/\cos A\)，得到

\[
\tan B=3\tan A
\]

\item 因为 \(\cos C=\frac{\sqrt{5}}{5}>0\)，所以 \(C\) 为锐角，且
\[
\sin C=\frac{2\sqrt{5}}{5}
\]
又 \(\tan C=2\). 由 \(A+B=\pi-C\)，所以 \(\tan(A+B)=-\tan C=-2\). 令 \(t=\tan A\)，由第（1）问得 \(\tan B=3t\)，从而
\[
\frac{4t}{1-3t^2}=-2
\]
由于 \(A+B=\pi-C\neq\frac{\pi}{2}\)，上式的分母不为零，整理得
\[
3t^2-2t-1=0
\]
即 \((3t+1)(t-1)=0\). 若 \(t=-\frac{1}{3}\)，则 \(A\)，\(B\) 均为钝角，不能构成三角形；故 \(t=1\). 因此 \(A=\frac{\pi}{4}\).
\end{enumerate}
\end{solution}

\begin{problem}
如图，在直三棱柱 \(ABC-A_{1}B_{1}C_{1}\) 中，\(A_{1}B_{1}=A_{1}C_{1}\)，\(D\)，\(E\) 分别是棱 \(BC\)，\(CC_{1}\) 上的点（点 \(D\) 不同于点 \(C\)），且 \(AD \perp DE\)，\(F\) 为 \(B_{1}C_{1}\) 的中点. 求证：

\begin{enumerate}
\item 平面 \(ADE \perp\) 平面 \(BCC_{1}B_{1}\)；
\item 直线 \(A_{1}F \parallel\) 平面 \(ADE\).
\end{enumerate}
\bitmapfigure[max width=0.42\linewidth]{img/2012/jiangsu/q16_fig1.png}
\end{problem}

\begin{solution}
由直三棱柱的性质，\(AB=A_{1}B_{1}\)，\(AC=A_{1}C_{1}\)，所以 \(AB=AC\). 因而 \(\triangle ABC\) 是等腰三角形.

因为 \(CC_{1}\perp\) 平面 \(ABC\)，而 \(E\) 在 \(CC_{1}\) 上，所以 \(DE\) 在底面上的射影是 \(DC\). 将 \(\overrightarrow{DE}\) 分解为 \(\overrightarrow{DC}+\overrightarrow{CE}\)，其中 \(\overrightarrow{CE}\perp\) 平面 \(ABC\)，而 \(\overrightarrow{AD}\) 在底面内，故
\[
\overrightarrow{AD}\cdot\overrightarrow{DE}=\overrightarrow{AD}\cdot\overrightarrow{DC}
\]
由已知 \(AD\perp DE\)，得 \(\overrightarrow{AD}\cdot\overrightarrow{DC}=0\).
由于 \(D\neq C\)，直线 \(DC\) 就是直线 \(BC\)，所以 \(AD\perp BC\). 在等腰三角形 \(ABC\) 中，顶点 \(A\) 到底边 \(BC\) 的垂线也是中线，故 \(D\) 是 \(BC\) 的中点.

\begin{enumerate}
\item 过 \(D\) 作直线 \(\ell\parallel CC_{1}\). 由于 \(D\in BC\)，且 \(BC\) 与 \(CC_{1}\) 都在平面 \(BCC_{1}B_{1}\) 内，所以 \(\ell\subset\) 平面 \(BCC_{1}B_{1}\). 又因为 \(AD\) 在平面 \(ABC\) 内，且 \(CC_{1}\perp\) 平面 \(ABC\)，所以 \(AD\perp\ell\). 已知 \(AD\perp BC\)，而 \(BC\) 与 \(\ell\) 相交于 \(D\)，故 \(AD\perp\) 平面 \(BCC_{1}B_{1}\). 由于 \(AD\subset\) 平面 \(ADE\)，可得平面 \(ADE\perp\) 平面 \(BCC_{1}B_{1}\).

\item \(D\) 是 \(BC\) 的中点，平移到上底面后，\(F\) 是 \(B_{1}C_{1}\) 的中点. 因而从 \(A_{1}\) 指向 \(F\) 的向量与从 \(A\) 指向 \(D\) 的向量平行，即 \(A_{1}F\parallel AD\). 又 \(DF\parallel CC_{1}\). 若 \(F\in DE\)，则 \(DE=DF\parallel CC_{1}\)，但 \(E\in DE\cap CC_{1}\)，矛盾；故 \(F\notin DE\). 若 \(A_{1}F\subset\) 平面 \(ADE\)，则 \(F\) 同时属于平面 \(ADE\) 和平面 \(BCC_{1}B_{1}\). 两平面的交线为 \(DE\)，这将推出 \(F\in DE\)，矛盾. 因此 \(A_{1}F\) 不在平面 \(ADE\) 内，且 \(A_{1}F\parallel AD\)、\(AD\subset\) 平面 \(ADE\)，所以直线 \(A_{1}F\parallel\) 平面 \(ADE\).
\end{enumerate}
\end{solution}

\begin{problem}
如图，建立平面直角坐标系 \(xOy\)，\(x\) 轴在地平面上，\(y\) 轴垂直于地平面，单位长度为 \(1\) 千米，某炮位于坐标原点. 已知炮弹发射后的轨迹在方程 \(y = kx - \frac{1}{20}(1+k^2)x^2\)（\(k>0\)）表示的曲线上，其中 \(k\) 与发射方向有关. 炮的射程是指炮弹落地点的横坐标.

\begin{enumerate}
\item 求炮的最大射程；
\item 设在第一象限有一飞行物（忽略其大小），其飞行高度为 \(3.2\) 千米，试问它的横坐标 \(a\) 不超过多少时，炮弹可以击中它？请说明理由.
\end{enumerate}
\bitmapfigure[max width=0.42\linewidth]{img/2012/jiangsu/q17_fig1.png}
\end{problem}

\begin{solution}
\begin{enumerate}
\item 炮弹落地时 \(y=0\). 除去发射点 \(x=0\)，由
\[
0=kx-\frac{1+k^2}{20}x^2
\]
得射程
\[
R(k)=\frac{20k}{1+k^2}=10-\frac{10(k-1)^2}{1+k^2}\leqslant10
\]
当且仅当 \(k=1\) 时等号成立，所以炮的最大射程为 \(10\) 千米.

\item 设飞行物的横坐标为 \(a\). 因为它在第一象限，所以 \(a>0\). 对固定的 \(a\)，轨迹在横坐标 \(a\) 处的高度为
\[
y(k)=-\frac{a^2}{20}k^2+ak-\frac{a^2}{20}=5-\frac{a^2}{20}-\frac{a^2}{20}\left(k-\frac{10}{a}\right)^2
\]
因此在 \(k>0\) 中，最大高度为 \(5-\frac{a^2}{20}\)，且取得最大值时 \(k=\frac{10}{a}>0\). 对固定的 \(a>0\)，上式右端关于 \(k\) 连续；将其在 \(k=0\) 处的值看作 \(-a^2/20<0\)，而在顶点处达到最大值. 所以由介值定理，存在 \(k>0\) 使炮弹经过飞行物的充要条件是
\[
5-\frac{a^2}{20}\geqslant3.2=\frac{16}{5}
\]
整理得 \(a^2\leqslant36\). 结合 \(a>0\)，可知 \(a\leqslant6\). 所以横坐标 \(a\) 不超过 \(6\) 千米时，炮弹可以击中飞行物.
\end{enumerate}
\end{solution}

\begin{problem}
若函数 \(y=f(x)\) 在 \(x=x_{0}\) 处取得极大值或极小值，则称 \(x_{0}\) 为函数 \(y=f(x)\) 的极值点. 已知 \(a\)，\(b\) 是实数，且 \(1\) 和 \(-1\) 是函数 \(f(x)=x^3+ax^2+bx\) 的两个极值点.

\begin{enumerate}
\item 求 \(a\) 和 \(b\) 的值；
\item 设函数 \(g(x)\) 的导函数 \(g'(x)=f(x)+2\)，求 \(g(x)\) 的极值点；
\item 设 \(h(x)=f(f(x))-c\)，其中 \(c\in[-2,2]\)，求函数 \(y=h(x)\) 的零点个数.
\end{enumerate}
\end{problem}

\begin{solution}
由 \(1\) 和 \(-1\) 是极值点，得导函数在这两个点处为零. 因此
\[
f'(x)=3x^2+2ax+b=3(x-1)(x+1)=3x^2-3
\]
比较系数可得 \(a=0\)，\(b=-3\). 此时 \(f'(x)=3(x^2-1)\) 在 \(-1\)，\(1\) 两侧变号，确实对应两个极值点.

\begin{enumerate}
\item \(a=0\)，\(b=-3\).

\item 此时 \(f(x)=x^3-3x\)，所以
\[
g'(x)=f(x)+2=x^3-3x+2=(x-1)^2(x+2)
\]
当 \(x<-2\) 时，\(g'(x)<0\)；当 \(-2<x<1\) 或 \(x>1\) 时，\(g'(x)>0\). 因此 \(g(x)\) 在 \(x=-2\) 的左侧递减、右侧递增，\(-2\) 是极小值点；虽然 \(g'(1)=0\)，但导函数在 \(1\) 两侧不变号，所以 \(1\) 不是极值点. 故 \(g(x)\) 的极值点为 \(-2\).

\item 令 \(u=f(x)\). 函数 \(f(u)=u^3-3u\) 的导数为 \(f'(u)=3(u^2-1)\)，所以它在 \((-\infty,-1)\)、\((1,+\infty)\) 上递增，在 \((-1,1)\) 上递减，并在 \(u=-1\) 处取得极大值 \(2\)，在 \(u=1\) 处取得极小值 \(-2\). 当 \(-2<c<2\) 时，方程
\[
f(u)=c
\]
在区间 \((-2,-1)\)、\((-1,1)\)、\((1,2)\) 内各有一个根：这是因为 \(f(-2)=-2\)，\(f(-1)=2\)，\(f(1)=-2\)，\(f(2)=2\)，再结合上述单调性. 对其中每一个根 \(u\in(-2,2)\)，方程 \(f(x)=u\) 也在这三个区间内各有一个根，且三根互不相同，因此 \(f(f(x))=c\) 有 \(3\times3=9\) 个零点.

当 \(c=2\) 时，
\[
f(u)-2=u^3-3u-2=(u-2)(u+1)^2
\]
所以 \(u\) 有两个不同的取值 \(2\)，\(-1\). 方程 \(f(x)=2\) 有两个不同的实根，方程 \(f(x)=-1\) 有三个不同的实根，共有 \(5\) 个零点. 当 \(c=-2\) 时，
\[
f(u)+2=u^3-3u+2=(u+2)(u-1)^2
\]
所以 \(u\) 有两个不同的取值 \(-2\)，\(1\). 方程 \(f(x)=-2\) 有两个不同的实根，且因为 \(1\in(-2,2)\)，由前面的单调性可知方程 \(f(x)=1\) 有三个不同的实根，共有 \(2+3=5\) 个零点. 故零点个数为
\[
\begin{cases}
9,&-2<c<2,\\
5,&c=-2\text{ 或 }c=2.
\end{cases}
\]
\end{enumerate}
\end{solution}

\begin{problem}
如图，在平面直角坐标系 \(xOy\) 中，椭圆 \(\frac{x^2}{a^2} + \frac{y^2}{b^2} = 1\)（\(a>b>0\)）的左、右焦点分别为 \(F_{1}(-c,0)\)，\(F_{2}(c,0)\). 已知点 \((1,e)\) 和 \(\left(e,\frac{\sqrt{3}}{2}\right)\) 都在椭圆上，其中 \(e\) 为椭圆的离心率.

\begin{enumerate}
\item 求椭圆的方程；
\item 设 \(A\)，\(B\) 是椭圆上位于 \(x\) 轴上方的两点，且直线 \(AF_{1}\) 与直线 \(BF_{2}\) 平行，\(AF_{2}\) 与 \(BF_{1}\) 交于点 \(P\).
\begin{enumerate}
\item 若 \(AF_{1}-BF_{2}=\frac{\sqrt{6}}{2}\)，求直线 \(AF_{1}\) 的斜率；
\item 求证：\(PF_{1}+PF_{2}\) 是定值.
\end{enumerate}
\end{enumerate}

\bitmapfigure[max width=0.42\linewidth]{img/2012/jiangsu/q19_fig1.png}
\end{problem}

\begin{solution}
\begin{enumerate}
\item 由离心率的定义，\(e^2=1-\frac{b^2}{a^2}\)，即 \(e^2=\frac{a^2-b^2}{a^2}\). 将点 \((1,e)\) 代入椭圆方程，得
\[
\frac{1}{a^2}+\frac{e^2}{b^2}=\frac{1}{a^2}+\frac{a^2-b^2}{a^2b^2}=\frac{1}{b^2}=1
\]
所以 \(b^2=1\). 再将点 \(\left(e,\frac{\sqrt{3}}{2}\right)\) 代入，得
\[
\frac{e^2}{a^2}+\frac{3}{4b^2}=1
\]
即
\[
\frac{a^2-1}{a^4}+\frac{3}{4}=1
\]
因此 \(a^4-4a^2+4=0\)，从而 \(a^2=2\). 故椭圆方程为
\[
\frac{x^2}{2}+y^2=1
\]

\item 由 \(a^2=2\)，\(b^2=1\)，得 \(c=1\)，所以 \(F_{1}=(-1,0)\)，\(F_{2}=(1,0)\). 若直线 \(AF_{1}\) 和 \(BF_{2}\) 都为竖直线，则 \(A=(-1,1/\sqrt{2})\)，\(B=(1,1/\sqrt{2})\). 此时两直线 \(AF_{2}\) 和 \(BF_{1}\) 的交点为 \(P=(0,1/(2\sqrt{2}))\)，从而
\[
PF_{1}+PF_{2}=2\sqrt{1+\frac{1}{8}}=\frac{3\sqrt{2}}{2}
\]
以下设两直线不为竖直线，记它们的斜率为 \(k\). 令
\[
A=(-1+t,kt)
\]
且 \(B=(1+s,ks)\)，
因为 \(A\)，\(B\) 位于 \(x\) 轴上方，所以 \(t\)，\(s\) 同号，且与 \(k\) 同号. 代入椭圆方程，得
\[
\begin{cases}
(1+2k^2)t^2-2t-1=0,\\
(1+2k^2)s^2+2s-1=0
\end{cases}
\]
记 \(d=1+2k^2\)，\(r=\sqrt{2(1+k^2)}\)，则 \(r>1\). 两个方程的根分别为
\(t=\frac{1\pm r}{d}\)，\(s=\frac{-1\pm r}{d}\).
当 \(k>0\) 时，\(t\)，\(s>0\)，故取正根；当 \(k<0\) 时，\(t\)，\(s<0\)，故取负根. 记 \(T=|t|\)，\(S=|s|\)，于是
\[
(T,S)=
\begin{cases}
\left(\dfrac{r+1}{d},\dfrac{r-1}{d}\right),& k>0,\\
\left(\dfrac{r-1}{d},\dfrac{r+1}{d}\right),& k<0.
\end{cases}
\]
因此无论 \(k\) 正负，都有
\[
TS=\frac{1}{d}
\]
且 \(S+T=\frac{2r}{d}\)，
并且 \(T>S\) 当且仅当 \(k>0\). 再记 \(q=\sqrt{1+k^2}\)，则
\[
AF_{1}=qT
\]
且 \(BF_{2}=qS\)，

\begin{enumerate}
\item 由 \(AF_{1}-BF_{2}=\frac{\sqrt{6}}{2}>0\)，可知 \(T>S\)，从而 \(k>0\). 此时
\[
T=\frac{1+r}{d}
\]
且 \(S=\frac{r-1}{d}\)，
从而 \(AF_{1}-BF_{2}=\frac{2q}{d}\)，
题设给出
\[
\frac{2q}{1+2k^2}=\frac{\sqrt{6}}{2}
\]
由于 \(1+2k^2=2q^2-1\)，平方并整理得
\[
12q^4-20q^2+3=0
\]
所以 \(q^2=\frac{3}{2}\) 或 \(q^2=\frac{1}{6}\). 因为 \(q^2=1+k^2\geqslant1\)，故 \(q^2=\frac{3}{2}\)，从而 \(k^2=\frac{1}{2}\). 又 \(k>0\)，所以直线 \(AF_{1}\) 的斜率为 \(\frac{\sqrt{2}}{2}\).

\item 椭圆上任意点到两焦点的距离和为 \(2a=2\sqrt{2}\)，故
\[
AF_{2}=2\sqrt{2}-qT
\]
且 \(BF_{1}=2\sqrt{2}-qS\)，
设 \(P=AF_{2}\cap BF_{1}\)，并写成
\[
P=F_{2}+\lambda(A-F_{2})=F_{1}+\mu(B-F_{1})
\]
由 \(A-F_{2}=(t-2,kt)\)，\(B-F_{1}=(s+2,ks)\)，比较横、纵坐标可得
\[
\lambda=\frac{S}{S+T}
\]
且 \(\mu=\frac{T}{S+T}\)，
因此
\[
PF_{2}=\frac{S}{S+T}AF_{2}
\]
且 \(PF_{1}=\frac{T}{S+T}BF_{1}\)，
于是
\[
PF_{1}+PF_{2}=\frac{T(2\sqrt{2}-qS)+S(2\sqrt{2}-qT)}{S+T}=2\sqrt{2}-\frac{2qTS}{S+T}=2\sqrt{2}-\frac{q}{r}=2\sqrt{2}-\frac{1}{\sqrt{2}}=\frac{3\sqrt{2}}{2}
\]
所以 \(PF_{1}+PF_{2}\) 是定值 \(\frac{3\sqrt{2}}{2}\).
\end{enumerate}
\end{enumerate}
\end{solution}

\begin{problem}
已知各项均为正数的两个数列 \(\{a_n\}\) 和 \(\{b_n\}\) 满足：\(a_{n+1} = \frac{a_n + b_n}{\sqrt{a_n^2 + b_n^2}}\)，\(n \in \N^*\).

\begin{enumerate}
\item 设 \(b_{n+1} = 1 + \frac{b_n}{a_n}\)，\(n \in \N^*\)，求证：数列 \(\left\{\left(\frac{b_n}{a_n}\right)^2\right\}\) 是等差数列；
\item 设 \(b_{n+1} = \sqrt{2} \cdot \frac{b_n}{a_n}\)，\(n \in \N^*\)，且 \(\{a_n\}\) 是等比数列，求 \(a_1\) 和 \(b_1\) 的值.
\end{enumerate}
\end{problem}

\begin{solution}
令 \(x_n=\frac{b_n}{a_n}>0\). 由题设递推式，
\[
a_{n+1}=\frac{1+x_n}{\sqrt{1+x_n^2}}
\]

\begin{enumerate}
\item 当 \(b_{n+1}=1+\frac{b_n}{a_n}=1+x_n\) 时，
\[
x_{n+1}=\frac{b_{n+1}}{a_{n+1}}=\frac{1+x_n}{(1+x_n)/\sqrt{1+x_n^2}}=\sqrt{1+x_n^2}
\]
所以 \(x_{n+1}^2=x_n^2+1\)，即数列 \(\{x_n^2\}=\left\{\left(\frac{b_n}{a_n}\right)^2\right\}\) 是以 \(1\) 为公差的等差数列.

\item 当 \(b_{n+1}=\sqrt{2}\,\frac{b_n}{a_n}=\sqrt{2}x_n\) 时，先设 \(\{a_n\}\) 的公比为 \(q>0\). 由上式及 \(x_n>0\)，有
\[
1<a_{n+1}=\frac{1+x_n}{\sqrt{1+x_n^2}}\leqslant\sqrt{2}
\]
因此对所有 \(n\geqslant2\)，都有 \(1<a_n\leqslant\sqrt{2}\). 若 \(q>1\)，则 \(a_n\) 趋于 \(+\infty\)；若 \(0<q<1\)，则 \(a_n\) 趋于 \(0\)，均与上述有界性矛盾. 所以 \(q=1\)，数列 \(\{a_n\}\) 为常数，记其值为 \(a\).

此时 \(b_{n+1}=\frac{\sqrt{2}}{a}b_n\)，所以 \(x_n=\frac{b_n}{a}\) 是以 \(r=\frac{\sqrt{2}}{a}\) 为公比的正数列. 又
\[
a=\frac{1+x_n}{\sqrt{1+x_n^2}}
\]
对所有 \(n\) 成立. 若 \(r>1\)，则 \(x_n\) 趋于 \(+\infty\)；若 \(0<r<1\)，则 \(x_n\) 趋于 \(0\). 这两种情况下等式右端都趋于 \(1\)，从而得到 \(a=1\)，与 \(a>1\) 矛盾. 因而 \(r=1\)，即 \(a=\sqrt{2}\).

于是
\[
\sqrt{2}=\frac{1+x_n}{\sqrt{1+x_n^2}}
\]
两边平方并整理，得 \((x_n-1)^2=0\)，所以 \(x_n=1\). 特别地，\(a_1=\sqrt{2}\)，\(b_1=a_1x_1=\sqrt{2}\).
\end{enumerate}
\end{solution}

\section{附加题：解答题}

\begin{problem}[A]
如图，\(AB\) 是圆 \(O\) 的直径，\(D\)，\(E\) 为圆 \(O\) 上位于 \(AB\) 异侧的两点，连接 \(BD\) 并延长至点 \(C\)，使 \(BD=DC\)，连接 \(AC\)，\(AE\)，\(DE\). 求证：\(\angle E=\angle C\).

\bitmapfigure[max width=0.42\linewidth]{img/2012/jiangsu/q21_fig1.png}
\end{problem}

\begin{solution}
因为 \(AB\) 是圆的直径，所以 \(\angle ADB=90^\circ\). 又 \(B\)，\(D\)，\(C\) 三点共线，故 \(AD\perp BC\). 由 \(BD=DC\)，点 \(D\) 是 \(BC\) 的中点. 因此在直角三角形 \(ADB\) 和 \(ADC\) 中，\(AD\) 为公共边且 \(BD=DC\)，两三角形全等，从而 \(AB=AC\). 所以 \(ABC\) 是等腰三角形，且 \(AD\) 同时是中线和高，有
\[
\angle ABD=\angle ABC=\angle BCA=\angle ACD
\]

点 \(A\)，\(B\)，\(D\)，\(E\) 在同一圆上，且由 \(D\)，\(E\) 位于 \(AB\) 异侧可知点 \(B\)，\(E\) 位于弦 \(AD\) 的同侧. 因此圆周角 \(\angle AED\) 与 \(\angle ABD\) 同对弦 \(AD\)，有 \(\angle AED=\angle ABD\). 因而
\[
\angle E=\angle AED=\angle ACD=\angle C
\]
\end{solution}

\reuseproblemnumber
\begin{problem}[B]
已知矩阵 \(A\) 的逆矩阵 \(A^{-1}=\begin{bmatrix}-\frac{1}{4}&\frac{3}{4}\\ \frac{1}{2}&-\frac{1}{2}\end{bmatrix}\)，求矩阵 \(A\) 的特征值.
\end{problem}

\begin{solution}
设 \(\mu\) 是矩阵 \(A^{-1}\) 的特征值，则
\[
0=\det\left(A^{-1}-\mu I\right)=\left(-\frac14-\mu\right)\left(-\frac12-\mu\right)-\frac38=\mu^2+\frac34\mu-\frac14
\]
解得 \(\mu=\frac14\) 或 \(\mu=-1\). 由于可逆矩阵的特征值与其逆矩阵的特征值互为倒数，矩阵 \(A\) 的特征值为 \(4\)，\(-1\).
\end{solution}

\reuseproblemnumber
\begin{problem}[C]
在极坐标系中，已知圆 \(C\) 经过点 \(P\left(\sqrt{2},\frac{\pi}{4}\right)\)，圆心为直线 \(\rho\sin\left(\theta-\frac{\pi}{3}\right)=-\frac{\sqrt{3}}{2}\) 与极轴的交点，求圆 \(C\) 的极坐标方程.
\end{problem}

\begin{solution}
圆心在极轴上，所以令 \(\theta=0\) 代入圆心所在直线，得
\[
\rho\sin\left(-\frac{\pi}{3}\right)=-\frac{\sqrt{3}}{2}
\]
从而 \(\rho=1\)，
因此圆心的直角坐标为 \((1,0)\). 点 \(P\left(\sqrt{2},\frac{\pi}{4}\right)\) 的直角坐标为 \((1,1)\)，所以圆的半径为 \(1\)，其直角坐标方程为
\[
(x-1)^2+y^2=1
\]
即 \(x^2+y^2-2x=0\). 代入 \(x=\rho\cos\theta\)，\(x^2+y^2=\rho^2\)，得圆的极坐标方程
\[
\rho^2-2\rho\cos\theta=0
\]
或写成 \(\rho=2\cos\theta\).
\end{solution}

\reuseproblemnumber
\begin{problem}[D]
已知实数 \(x\)，\(y\) 满足：\(|x+y|<\frac{1}{3}\)，\(|2x-y|<\frac{1}{6}\)，求证：\(|y|<\frac{5}{18}\).
\end{problem}

\begin{solution}
令 \(u=x+y\)，\(v=2x-y\). 由
\[
2u-v=2(x+y)-(2x-y)=3y
\]
故
\[
y=\frac{2u-v}{3}
\]
又因 \(|u|<\frac13\)、\(|v|<\frac16\)，由绝对值三角不等式，
\[
|y|=\frac{|2u-v|}{3}\leqslant\frac{2|u|+|v|}{3}<\frac{2\cdot\frac13+\frac16}{3}=\frac{5}{18}
\]
所以 \(|y|<\frac{5}{18}\).
\end{solution}

\begin{problem}
设 \(\xi\) 为随机变量，从棱长为 \(1\) 的正方体的 \(12\) 条棱中任取两条，当两条棱相交时，\(\xi=0\)；当两条棱平行时，\(\xi\) 的值为两条棱之间的距离；当两条棱异面时，\(\xi=1\).

\begin{enumerate}
\item 求概率 \(P(\xi=0)\)；
\item 求 \(\xi\) 的分布列，并求其数学期望 \(E(\xi)\).
\end{enumerate}
\end{problem}

\begin{solution}
从 \(12\) 条棱中任取两条，共有
\[
\mathrm C_{12}^{2}=66
\]
种取法.

\begin{enumerate}
\item 正方体的每个顶点有 \(3\) 条棱相交于该点，因此相交的棱对有
\[
8\mathrm C_{3}^{2}=24
\]
对. 每一对相交棱只有一个公共顶点，所以没有重复计数. 因而
\[
P(\xi=0)=\frac{24}{66}=\frac{4}{11}
\]

\item 按棱的方向分类，正方体有 \(3\) 个方向，每个方向有 \(4\) 条互相平行的棱，所以平行的棱对有
\[
3\mathrm C_{4}^{2}=18
\]
对. 对于同一方向的 \(4\) 条棱，在垂直于该方向的正方形中对应四个顶点；其中相邻顶点给出距离 \(1\)，有 \(4\) 对，两个对角顶点给出距离 \(\sqrt{2}\)，有 \(2\) 对. 因此
\[
N(\xi=1\text{ 且两棱平行})=3\times4=12
\]
且 \(N(\xi=\sqrt{2})=3\times2=6\)，

其余棱对为异面棱对，数目为
\[
66-24-18=24
\]
这些棱对也满足 \(\xi=1\). 所以 \(P(\xi=0)=\frac{4}{11}\)，\(P(\xi=1)=\frac{6}{11}\)，\(P(\xi=\sqrt{2})=\frac{1}{11}\)，这就是 \(\xi\) 的分布列.
从而
\[
E(\xi)=0\cdot\frac{4}{11}+1\cdot\frac{6}{11}+\sqrt{2}\cdot\frac{1}{11}=\frac{6+\sqrt{2}}{11}
\]
\end{enumerate}
\end{solution}

\begin{problem}
设集合 \(P_n = \{1,2,\cdots,n\}\)，\(n \in \N^*\). 记 \(f(n)\) 为同时满足下列条件的集合 \(A\) 的个数：
\begin{circlelist}
\item \(A \subseteq P_n\)；
\item 若 \(x \in A\)，则 \(2x \notin A\)；
\item 若 \(x \in C_{P_n}A\)，则 \(2x \notin C_{P_n}A\).
\end{circlelist}

\begin{enumerate}
\item 求 \(f(4)\)；
\item 求 \(f(n)\) 的解析式（用 \(n\) 表示）.
\end{enumerate}
\end{problem}

\begin{solution}
对每个 \(j\in P_n\)，都可以唯一地写成 \(j=2^r u\)，其中 \(u\) 为奇数. 因而 \(P_n\) 被分成若干条链
\(u\)，\(2u\)，\(4u\)，\(\ldots\)，\(2^m u\leqslant n\)
其中每条链的首项都是不超过 \(n\) 的奇数.

在同一条链上，除最后一项外，每一项的后继仍属于 \(P_n\). 对这样的相邻两项，条件 2 表明若前一项属于 \(A\)，则后一项不属于 \(A\)；条件 3 表明若前一项不属于 \(A\)，则后一项属于 \(A\). 因此同一条链上的成员资格必须交替出现，而只要确定首项 \(u\) 是否属于 \(A\)，整条链的取法就唯一确定. 不同链之间互不影响.

不超过 \(n\) 的奇数共有 \(\left\lceil\frac{n}{2}\right\rceil\) 个，所以
\[
f(n)=2^{\left\lceil\frac{n}{2}\right\rceil}
\]
当 \(n=4\) 时，奇数首项为 \(1\)，\(3\)，故 \(f(4)=2^2=4\).
\end{solution}
